% docs/slides/preamble.tex -- 五份 deck 共享
\documentclass[aspectratio=169,11pt]{beamer}

% ==== 中文与字体 ====
\usepackage[UTF8,fontset=none]{ctex}
\usepackage{fontspec}

\IfFontExistsTF{PingFang SC}{%
  \setCJKmainfont{PingFang SC}[BoldFont=PingFang SC,ItalicFont=PingFang SC,AutoFakeSlant=0.15]
  \setCJKsansfont{PingFang SC}[ItalicFont=PingFang SC,AutoFakeSlant=0.15]
  \setCJKmonofont{PingFang SC}
}{%
  \IfFontExistsTF{Source Han Sans SC}{%
    \setCJKmainfont{Source Han Sans SC}
    \setCJKsansfont{Source Han Sans SC}
  }{%
    \IfFontExistsTF{Noto Sans CJK SC}{%
      \setCJKmainfont{Noto Sans CJK SC}
      \setCJKsansfont{Noto Sans CJK SC}
    }{%
      \setCJKmainfont{Heiti SC}
      \setCJKsansfont{Heiti SC}
    }
  }
}

\IfFontExistsTF{Helvetica Neue}{\setmainfont{Helvetica Neue}}{%
  \IfFontExistsTF{Inter}{\setmainfont{Inter}}{}}
\IfFontExistsTF{JetBrains Mono}{\setmonofont{JetBrains Mono}[Scale=0.85]}%
  {\setmonofont{Menlo}[Scale=0.85]}

% ==== 颜色 ====
\usepackage{xcolor}
\definecolor{cMain}{HTML}{1F4E79}
\definecolor{cAccent}{HTML}{E07A1F}
\definecolor{cGray}{HTML}{595959}
\definecolor{cMute}{HTML}{F4F4F4}
\definecolor{cOK}{HTML}{2E7D32}
\definecolor{cBad}{HTML}{B23A3A}
\definecolor{cWarn}{HTML}{C8A300}

% ==== 主题 ====
\usepackage{tcolorbox}
\usetheme{metropolis}
\metroset{progressbar=frametitle,numbering=fraction,block=fill}
\setbeamercolor{frametitle}{bg=cMain,fg=white}
\setbeamercolor{progress bar}{fg=cAccent,bg=cMute}
\setbeamercolor{alerted text}{fg=cAccent}
\setbeamercolor{title}{fg=cMain}

% metropolis 默认尝试 Fira Sans，未安装时 fallback 到 Latin Modern Sans。
% 这里在主题加载后显式覆盖，强制使用 macOS 上一定存在的 Helvetica Neue，
% 跟中文 PingFang 风格一致；找不到则按顺序尝试 Inter / Arial。
\IfFontExistsTF{Helvetica Neue}{%
  \setsansfont{Helvetica Neue}%
}{%
  \IfFontExistsTF{Inter}{\setsansfont{Inter}}{%
    \IfFontExistsTF{Arial}{\setsansfont{Arial}}{}}%
}

% ==== TikZ ====
\usepackage{tikz}
\usetikzlibrary{positioning,arrows.meta,fit,shapes.geometric,calc,%
  decorations.pathmorphing,backgrounds}

\tikzset{
  node-box/.style={draw=cMain,thick,rounded corners=2pt,
    minimum height=8mm,minimum width=18mm,align=center,font=\small},
  node-state/.style={draw=cMain,thick,circle,minimum size=14mm,
    align=center,font=\small},
  node-ok/.style={draw=cOK,thick,fill=cOK!10},
  node-warn/.style={draw=cWarn,thick,fill=cWarn!15},
  node-bad/.style={draw=cBad,thick,fill=cBad!10},
  % 色盲友好的状态机三态：蓝 (正常/Closed) / 橙 (告警/Open) / 灰 (中间/HalfOpen)
  % 避免 red-green 二元区分，改用 hue + lightness 双通道。
  node-state-normal/.style={draw=cMain,thick,fill=cMain!12},
  node-state-alert/.style={draw=cAccent,thick,fill=cAccent!18},
  node-state-transition/.style={draw=cGray,thick,fill=cGray!18},
  arrow/.style={-{Stealth[length=2mm]},thick,draw=cMain},
  arrow-async/.style={-{Stealth[length=2mm]},thick,dashed,draw=cGray},
  arrow-bad/.style={-{Stealth[length=2mm]},thick,draw=cBad},
  arrow-alert/.style={-{Stealth[length=2mm]},thick,draw=cAccent},
  lifeline/.style={dashed,thick,draw=cGray},
}

% ==== 代码 ====
\usepackage{listings}
\lstdefinelanguage{Go}{
  morekeywords={break,case,chan,const,continue,default,defer,else,
    fallthrough,for,func,go,goto,if,import,interface,map,package,range,
    return,select,struct,switch,type,var,nil,true,false,iota},
  morekeywords=[2]{string,int,int64,uint,uint64,bool,byte,error,context},
  sensitive=true,
  morecomment=[l]//,morecomment=[s]{/*}{*/},
  morestring=[b]",morestring=[b]`,
}
\lstdefinelanguage{LuaRedis}{
  morekeywords={and,break,do,else,elseif,end,false,for,function,if,in,
    local,nil,not,or,repeat,return,then,true,until,while},
  morekeywords=[2]{redis,call,KEYS,ARGV,tonumber,HGET,HSET,GET,SET,
    DECR,DECRBY,INCR,INCRBY,EXPIRE,PEXPIRE,ZADD,ZCARD,ZREMRANGEBYSCORE,
    EVAL,EVALSHA},
  sensitive=true,
  morecomment=[l]{--},
  morestring=[b]",morestring=[b]',
}

\lstdefinestyle{base}{
  basicstyle=\ttfamily\scriptsize,
  numbers=left,numberstyle=\tiny\color{cGray},
  keywordstyle=\color{cMain}\bfseries,
  keywordstyle=[2]\color{cAccent},
  stringstyle=\color{cOK},
  commentstyle=\color{cGray}\itshape,
  backgroundcolor=\color{cMute},
  frame=single,framerule=0pt,
  showstringspaces=false,
  breaklines=true,
  tabsize=2,
  columns=fullflexible,
}
\lstset{style=base}

% ==== 三线表与附录页码 ====
\usepackage{booktabs}
\usepackage{appendixnumberbeamer}

% ==== 页脚来源标注 ====
\newcommand{\srcnote}[1]{{\color{cGray}\scriptsize\itshape #1}}

% 关键论点框
\newcommand{\keypoint}[1]{%
  \begin{tcolorbox}[colback=cMute,colframe=cMain,boxrule=0.5pt,arc=2pt]%
  #1\end{tcolorbox}}


\title{30 分钟那道线}
\subtitle{关单时机的业务博弈与双保险机制 (gomall 业务侧 deck)}
\author{gomall · 业务侧 deck}
\date{}

\begin{document}

\maketitle

\begin{frame}{今日大纲}
\tableofcontents
\end{frame}

% ============================================================
\section{30 分钟是怎么定的}
% ============================================================

\begin{frame}{这道线两头都得罪人}
\begin{tikzpicture}[node distance=0mm and 0mm,every node/.style={font=\scriptsize}]
  \draw[arrow] (0,0) -- (13,0);
  \foreach \x/\lbl in {0/0min,2/5min,4/15min,6.5/30min,9/60min,12/24h}{
    \draw[cGray] (\x,-0.08) -- (\x,0.08);
    \node[below=1mm] at (\x,-0.08) {\lbl};
  }
  \node[node-box,minimum width=18mm,fill=cBad!10,draw=cBad,above=4mm of {(2,0)}]
    (s1) {商家投诉区};
  \node[node-box,minimum width=22mm,fill=cOK!10,draw=cOK,above=4mm of {(6.5,0)}]
    (s2) {\textbf{业务平衡点}};
  \node[node-box,minimum width=22mm,fill=cBad!10,draw=cBad,above=4mm of {(11,0)}]
    (s3) {库存周转崩盘区};
\end{tikzpicture}
\vspace{4mm}
\begin{itemize}
  \item 商家侧抱怨："我刚下完单 5 分钟就被取消"——5 分钟切支付方式 / 等验证码都不够。
  \item 运营侧抱怨："双 11 锁单 24 小时谁也抢不到"——爆款被恶意囤一整天，正经买家流失。
  \item 30 分钟是\textbf{双方都能接受的最大公约数}，不是技术参数，是\textbf{业务的合同条款}。
\end{itemize}
\srcnote{repository/rabbitmq/order\_delay.go:22 OrderCancelDelay = 30 * time.Minute}
\end{frame}

\begin{frame}{为什么不是 5min / 15min / 60min}
\begin{tabular}{@{}llp{5cm}@{}}
\toprule
方案 & 用户体感 & 业务后果 \\
\midrule
5 分钟    & 切卡 / 等短信都不够 & 商家投诉"刚下单就被取消"，付款转化掉 10\%+ \\
15 分钟   & 多数人够用         & 高峰期仍有 5--8\% 误关，弱网用户感受不友好 \\
\textbf{30 分钟} & 含切卡 + 短信冗余  & 大促锁单可控，商家 / 运营都能接受 \\
60 分钟   & 用户不抱怨         & 双 11 锁 1 小时，爆款一直显示"已售罄" \\
24 小时   & 客服零投诉         & 库存周转崩盘，恶意囤货成本为零 \\
\bottomrule
\end{tabular}
\vspace{2mm}
\small 30min 是\textbf{工业界经验值}：京东 30min / 天猫 30min--1h / 拼多多 15min。改这个值要谈\textbf{5 个角色}：商家（锁单投诉）/ 运营（大促临时缩短）/ 风控（黑产窗口）/ 客服（付款重试）/ 财务（待支付科目）——代码侧只是把决策落到 \texttt{OrderCancelDelay} 一行。
\srcnote{repository/rabbitmq/order\_delay.go:21--22 30min 常量}
\end{frame}

% ============================================================
\section{短 / 长的代价：双 11 的两个真实场景}
% ============================================================

\begin{frame}{场景 A：5 分钟太短 = 付款流失}
\begin{tikzpicture}[node distance=8mm and 12mm]
  \node[node-box] (order) {下单成功\\\scriptsize T = 0};
  \node[node-box,right=of order,node-warn] (pay1) {切到银行 App\\\scriptsize +90s};
  \node[node-box,right=of pay1,node-warn] (sms) {等短信验证码\\\scriptsize +120s};
  \node[node-box,right=of sms,node-warn] (slow) {弱网下载页面\\\scriptsize +60s};
  \node[node-box,below=of sms,node-bad] (cancel) {订单被取消\\\scriptsize T = 5min};
  \draw[arrow] (order) -- (pay1);
  \draw[arrow] (pay1) -- (sms);
  \draw[arrow] (sms) -- (slow);
  \draw[arrow-bad] (slow) -- (cancel);
\end{tikzpicture}
\vspace{2mm}
\begin{itemize}
  \item 5 分钟里要做：切银行 / 等短信 / 弱网下载 / 输密码 / 人脸识别。
  \item 客单价越高，付款链路越长（双卡 / 分期 / 信用卡确认）。
  \item \textbf{真实数字}：电商支付链路 p95 = 120s，p99 = 240s。5min 砍 p99 用户。
\end{itemize}
\srcnote{支付链路时延来自业界 RUM 调研，非 gomall 实测}
\end{frame}

\begin{frame}{场景 B：24 小时太长 = 恶意囤货}
\begin{tikzpicture}[node distance=6mm and 10mm]
  \node[node-box,node-bad] (bot) {黑产脚本\\\scriptsize 1000 账号};
  \node[node-box,right=of bot,node-warn] (lock) {批量下单\\\scriptsize 每号锁 5 件};
  \node[node-box,right=of lock,node-bad] (out) {商品全部"已售罄"\\\scriptsize 5000 件};
  \node[node-box,below=of lock,node-warn] (real) {正经用户\\\scriptsize 看到"无货"};
  \node[node-box,below=of out,node-bad] (gone) {流失};
  \draw[arrow] (bot) -- (lock);
  \draw[arrow-bad] (lock) -- (out);
  \draw[arrow] (out) -- (real);
  \draw[arrow-bad] (real) -- (gone);
\end{tikzpicture}
\vspace{2mm}
\begin{itemize}
  \item 24 小时锁单 = 黑产\textbf{零成本}囤光商品（不付款，第二天自动取消）。
  \item 商家库存看上去"卖完了"，运营找不到原因。
  \item 真实用户进来看到"已售罄"$\to$ 流失。GMV 没增长，库存周转崩盘。
  \item 这是\textbf{库存被武器化}：把"自动关单延迟"变成攻击工具。
\end{itemize}
\srcnote{deck 11 inventory-business 第 4 节 恶意锁单 同主题}
\end{frame}

% ============================================================
\section{RMQ + Cron 双保险：单链路都不够}
% ============================================================

\begin{frame}{为什么不能只留一条链路}
\begin{tikzpicture}[node distance=10mm and 14mm]
  \node[node-box] (order) {下单};
  \node[node-box,right=of order,node-warn] (mq) {RMQ 延迟队列\\\scriptsize TTL 30min};
  \node[node-box,right=of mq,node-ok] (close1) {关单};
  \node[node-box,below=of mq,node-warn] (cron) {Cron 兜底\\\scriptsize 扫 UnPaid + 15min};
  \draw[arrow] (order) -- (mq);
  \draw[arrow] (mq) -- (close1);
  \draw[arrow,dashed] (order) to[bend right=20] (cron);
  \draw[arrow,dashed] (cron) to[bend right=20] (close1);
\end{tikzpicture}
\vspace{2mm}
\begin{itemize}
  \item \textbf{只 RMQ}：进程挂 / 队列丢消息 / nack 进 DLX = 订单停在 UnPaid，库存永远不释放。
  \item \textbf{只 Cron}：5min 间隔 + 单批 100 条上限，库存释放 lag 偏大；高峰期跟不上下单速度。
  \item \textbf{双保险}：RMQ 主路径（30min 准点）+ Cron 兜底（15min 窗口 / 5min 节奏）。
  \item \textbf{为什么 Cron 是 15min 窗口而非 30min}：RMQ 正常时 Cron 扫到的订单到 30min 已被 RMQ 关掉（WHERE type=UnPaid 失败）→ Cron 不抢功；RMQ 挂掉时 Cron 立刻接管，库存最多多锁 15--20min。
\end{itemize}
\srcnote{repository/rabbitmq/order\_delay.go:24--56 拓扑；service/order\_task.go:14--48 + :16 GetTimeoutOrders(15,100)}
\end{frame}

\begin{frame}{RMQ 延迟队列拓扑：死信交换换出"延迟"}
\begin{tikzpicture}[node distance=6mm and 14mm,every node/.style={font=\small}]
  \node[node-box] (prod) {producer\\\scriptsize PublishOrderCancelDelay};
  \node[node-box,right=of prod,node-warn] (ex1) {order.delay\\.exchange};
  \node[node-box,right=of ex1,node-warn] (q1) {order.delay.queue\\\scriptsize TTL 30min};
  \node[node-box,below=of q1,node-warn] (ex2) {order.dead\\.exchange};
  \node[node-box,left=of ex2,node-warn] (q2) {order.dead.queue};
  \node[node-box,left=of q2,node-ok] (cons) {consumer\\\scriptsize CancelUnpaidOrder};
  \draw[arrow] (prod) -- (ex1);
  \draw[arrow] (ex1) -- (q1);
  \draw[arrow-async] (q1) -- node[right,font=\tiny]{TTL 过期} (ex2);
  \draw[arrow] (ex2) -- (q2);
  \draw[arrow] (q2) -- (cons);
\end{tikzpicture}
\vspace{2mm}
\begin{itemize}
  \item RabbitMQ 没有原生延迟队列——\textbf{TTL + 死信交换}是工业界标准做法。
  \item 消息进 delay 队列后 30min "未被消费"$\to$ TTL 触发 $\to$ 流向 dead 交换 $\to$ 真正的消费者吃。
  \item 持久化（\texttt{DeliveryMode: Persistent}） + 手动 ack，避免消费者重启丢消息。
\end{itemize}
\srcnote{repository/rabbitmq/order\_delay.go:24--56 拓扑；:58--72 publish；:74--105 consume}
\end{frame}


\begin{frame}{RMQ 启动失败时的静默降级 = 双保险的最佳实证}
\begin{tikzpicture}[node distance=8mm and 12mm]
  \node[node-box] (boot) {cmd/main.go};
  \node[node-box,right=of boot,node-warn] (try) {tryInitRabbitMQ};
  \node[node-box,right=of try,node-bad] (fail) {RMQ 未启动\\\scriptsize panic};
  \node[node-box,below=of try,node-warn] (rec) {defer recover\\\scriptsize Warnf};
  \node[node-box,right=of rec,node-ok] (cron) {Cron 仍然跑\\\scriptsize 兜底接管};
  \draw[arrow] (boot) -- (try);
  \draw[arrow-bad] (try) -- (fail);
  \draw[arrow,dashed] (fail) to[bend right=20] (rec);
  \draw[arrow] (rec) -- (cron);
\end{tikzpicture}
\vspace{2mm}
\begin{itemize}
  \item gomall 压测报告里写得清楚：\textbf{"RabbitMQ 未启动 (不影响 HTTP 链路)"}。
  \item RMQ 启动失败 = recover 捕获 panic + 打 Warn 日志，HTTP 服务继续接客。
  \item 关单链路降级为\textbf{只走 Cron}，业务承诺打折但不崩。
  \item 这是\textbf{双保险设计的最佳实证}：缺一不可，但缺一不死。
\end{itemize}
\srcnote{cmd/main.go:49--58 tryInitRabbitMQ 静默降级；stressTest/REPORT.md 第 7 行}
\end{frame}

% ============================================================
\section{CancelUnpaidOrder 必须幂等}
% ============================================================

\begin{frame}{至少 4 个调用方会触发同一笔关单}
\begin{tikzpicture}[node distance=6mm and 10mm]
  \node[node-box,node-warn] (mq)   {RMQ 消费者\\\scriptsize 30min TTL};
  \node[node-box,node-warn,right=of mq] (cron) {Cron 兜底\\\scriptsize 15min 扫};
  \node[node-box,node-warn,right=of cron] (cs) {客服后台\\\scriptsize 手动取消};
  \node[node-box,node-warn,right=of cs] (user) {用户主动\\\scriptsize APP 取消};
  \node[node-box,below=of cs,node-ok] (target) {CancelUnpaidOrder\\\scriptsize 同一笔订单};
  \draw[arrow] (mq)   -- (target);
  \draw[arrow] (cron) -- (target);
  \draw[arrow] (cs)   -- (target);
  \draw[arrow] (user) -- (target);
\end{tikzpicture}
\vspace{2mm}
\begin{itemize}
  \item 同一笔订单可能\textbf{同时}被 4 个来源触发关单（高并发下并不罕见）。
  \item 例：用户付款失败手动取消 → Cron 也扫到 → RMQ 也到期 → 三次调用打到 \texttt{CancelUnpaidOrder}。
  \item 如果不幂等：库存释放 3 次（实际欠库）、outbox 写 3 条 \texttt{order.cancelled}（下游处理 3 次）。
\end{itemize}
\srcnote{service/order\_cancel.go:19--63 单一入口接 4 个调用方}
\end{frame}

\begin{frame}[fragile]{CloseOrderWithCheck：用 WHERE 条件兜住幂等}
\begin{lstlisting}[language=Go,firstnumber=185,
  caption={repository/db/dao/order.go:185--196},captionpos=b]
func (dao *OrderDao) CloseOrderWithCheck(orderNum uint64) (bool, error) {
    res := dao.DB.Model(&model.Order{}).Where(
        "order_num=? and type=?", orderNum, consts.UnPaid).
        Update("type", consts.Cancelled)
    if res.Error != nil {
        return false, res.Error
    }
    return res.RowsAffected > 0, nil
}
\end{lstlisting}
\vspace{1mm}
\begin{itemize}
  \item \textbf{关键是 WHERE 里那个 \texttt{type=UnPaid}}：只有 UnPaid 才会被改成 Cancelled，已 Cancelled / Paid 直接 RowsAffected=0。
  \item 第二次调用看到 RowsAffected=0 → \texttt{closed = false} → 直接 return，不释放库存、不写 outbox。
  \item \textbf{一行 SQL 兜住了 4 个调用方的并发竞争}。
\end{itemize}
\srcnote{repository/db/dao/order.go:185--196 幂等关单 SQL；service/order\_cancel.go:30--57 closed 旗标}
\end{frame}

\begin{frame}{幂等失守会发生什么}
\begin{tabular}{@{}lll@{}}
\toprule
失守场景 & 用户侧后果 & 业务侧后果 \\
\midrule
库存释放 N 次 & 看到"虚高"库存 & 实际欠库，超卖必赔 \\
outbox 发 N 条 cancelled & 下游收 N 次 & ES 删 N 次、商家通知 N 次 \\
状态机回退 (Paid → Cancelled) & 钱付了订单消失 & 资损 P0 \\
\bottomrule
\end{tabular}
\vspace{2mm}
\begin{itemize}
  \item 最严重的是\textbf{状态机回退}：用户已经付了款的订单被误关单 = 资损。
  \item gomall 用 \texttt{type=UnPaid} 这个 WHERE 死锁住——已付款的订单永远不会被 cancel 路径碰到。
  \item 这条 SQL 是关单链路的\textbf{最后一道防线}，绝对不能改。
\end{itemize}
\srcnote{repository/db/dao/order.go:185--196；consts/order.go UnPaid / Paid / Cancelled 状态码}
\end{frame}

% ============================================================
\section{关单是金融事件，不是状态变更}
% ============================================================

\begin{frame}{一次关单 = 三个动作必须\textbf{原子}发生}
\begin{tikzpicture}[node distance=8mm and 12mm]
  \node[node-box,node-warn] (close) {1. CloseOrderWithCheck\\\scriptsize type UnPaid → Cancelled};
  \node[node-box,node-warn,right=of close] (outbox) {2. outbox 写\\\scriptsize order.cancelled};
  \node[node-box,below=of close,node-warn] (release) {3. release reservation\\\scriptsize Redis 桶};
  \node[node-box,below=of outbox,node-ok] (tx) {同一个 DB tx\\\scriptsize 包 1 + 2};
  \draw[arrow] (close) -- (outbox);
  \draw[arrow-async] (outbox) -- (release);
  \draw[arrow,dashed] (close) -- (tx);
  \draw[arrow,dashed] (outbox) -- (tx);
\end{tikzpicture}
\vspace{2mm}
\begin{itemize}
  \item \textbf{1 + 2 必须在同一个 DB 事务}：状态变更和事件发布原子绑定（outbox 模式）。
  \item \textbf{3 在 tx 之外}：Redis 不支持 2PC，failed → 留 reserved 给下次对账修正，胜过参与分布式事务。
  \item 这是\textbf{关单作为金融事件}的最低要求：状态 + 事件 + 库存三件套必须协调。
\end{itemize}
\srcnote{service/order\_cancel.go:31--51 DB tx 包 1+2；:59--61 tx 外释放 Redis}
\end{frame}

\begin{frame}{不能给商家发"发货通知"}
\begin{itemize}
  \item 关单事件 \texttt{order.cancelled} 的下游消费者绝对\textbf{不能}包含"通知商家发货"——
  \begin{itemize}
    \item 商家收到"发货"通知 → 真发货 → 用户没付款 → 商家亏运费 + 货损。
  \end{itemize}
  \item 关单事件的合理下游：\textbf{ES 删除索引、库存对账、客服工单关闭、消息推送"订单已取消"}。
  \item \textbf{关键不变量}：\texttt{order.cancelled} 永远只对应 type=UnPaid 的订单（已 Paid 走的是\texttt{order.refunded}）。
  \item 业务侧：客服收到"订单已取消"工单时，不需要联系商家——状态机保证。
\end{itemize}
\vspace{2mm}
\keypoint{关单不是简单的状态翻转，是触发\textbf{一组下游业务动作}的金融事件——\\
谁该消费、谁不该消费，必须在事件设计时就定死。}
\vspace{1mm}
\small \textbf{对账兜底}：凌晨低峰跑 \texttt{SUM(reserved)} vs \texttt{SUM(UnPaid.Num)}，偏差告警——双保险都漏时的最后一道闸。当前 gomall 运维 backlog 第 1 项，待落地。
\srcnote{service/order\_cancel.go:40--50 outbox payload；repository/cache/inventory.go:127--138 对账入口}
\end{frame}


% ============================================================
\section{遗留与改进}
% ============================================================

\begin{frame}{遗留 1：三处时长口径不一致}
\begin{tabular}{@{}lll@{}}
\toprule
位置 & 数值 & 用途 \\
\midrule
\texttt{repository/rabbitmq/order\_delay.go:22} & 30min & RMQ TTL（业务关单时机）\\
\texttt{service/order\_task.go:16} & 15min & Cron 扫表窗口（兜底安全余量）\\
\texttt{service/order.go:88--92} & 15min & Redis ZSet score 偏移 \\
\bottomrule
\end{tabular}
\vspace{2mm}
\begin{itemize}
  \item 三处口径\textbf{业务含义不同}（关单时机 / 兜底窗口 / 内部排序），但都涉及"15min 这个数"。
  \item 风险：未来改 30min → 15min 时，是否要同步改 Cron / ZSet？
  \item 改进：抽到 \texttt{consts/order.go} 单独命名（\texttt{OrderCloseTTL} / \texttt{OrderCronBacklog} / \texttt{OrderZSetOffset}）。
\end{itemize}
\srcnote{repository/rabbitmq/order\_delay.go:22；service/order\_task.go:16；service/order.go:88--92}
\end{frame}

\begin{frame}{遗留 2：\texttt{orders/delete} 绕过 CancelUnpaidOrder}
\begin{tikzpicture}[node distance=8mm and 12mm]
  \node[node-box,node-warn] (route) {POST /orders/delete\\\scriptsize router.go:81};
  \node[node-box,right=of route,node-bad] (svc) {OrderDelete};
  \node[node-box,right=of svc,node-bad] (sql) {DeleteOrderById\\\scriptsize 软删除 only};
  \node[node-box,below=of svc,node-ok] (correct) {应当：CancelUnpaidOrder};
  \draw[arrow-bad] (route) -- (svc);
  \draw[arrow-bad] (svc) -- (sql);
  \draw[arrow,dashed] (svc) -- (correct);
\end{tikzpicture}
\vspace{2mm}
\begin{itemize}
  \item \texttt{routes/router.go:81} 的 \texttt{orders/delete} 直接调 \texttt{OrderDelete} → \texttt{DeleteOrderById}（GORM 软删除）。
  \item \textbf{漏了 release reservation 和 outbox 写}——用户"删除"未支付订单后，Redis 桶仍占着库存，下游 ES 不知道订单消失。
  \item 这是一条\textbf{绕过关单链路}的捷径，建议要么对 UnPaid 改走 CancelUnpaidOrder，要么禁掉这个 endpoint。
\end{itemize}
\srcnote{routes/router.go:81；service/order.go:168--192 OrderDelete；repository/db/dao/order.go:149--154 DeleteOrderById}
\end{frame}

\begin{frame}[fragile]{遗留 3：Cron 表达式 bug——意图 5min 一次, 实际 60 次}
\begin{lstlisting}[language=Go,firstnumber=14,
  caption={initialize/cron.go:14--23 当前代码},captionpos=b]
c := cron.New(cron.WithSeconds())
orderService := new(service.OrderTaskService)
_, err := c.AddFunc("* */5 * * * *", func() {
    defer func() {
        if r := recover(); r != nil { /* ... */ }
    }()
    orderService.RunOrderTimeoutCheck()
})
\end{lstlisting}
\vspace{1mm}
\begin{itemize}
  \item 表达式 \texttt{* */5 * * * *} 含义是 "每月 / 每周日 / 每天 / 每 5 分钟 / 每分钟内每秒"——\textbf{即 5min 里每秒都触发 60 次}。
  \item 业务意图是 "每 5 分钟跑一次"，正确写法是 \texttt{0 */5 * * * *}（秒位 = 0）。
  \item 影响：Cron 调度器 5min 内调用 RunOrderTimeoutCheck 60 次，\textbf{DB 被无谓打了 60 倍}。
  \item 幸运的是 CloseOrderWithCheck 幂等，没有功能性 bug——但成本暴涨且日志告警淹没。
\end{itemize}
\srcnote{initialize/cron.go:16；robfig/cron/v3 with-seconds 字段顺序}
\end{frame}

\begin{frame}{遗留 4：RMQ 长时间不通 / Cron 长时间不跑都没告警}
\begin{itemize}
  \item RMQ 静默降级\textbf{很优雅}（cmd/main.go:49--58 recover），但运维侧\textbf{看不到}——只有一行 Warn 日志，没有 metric。
  \item Cron 进程如果死掉（panic 没 recover / 进程被 OOM kill），\textbf{没有外部心跳监控}知道。
  \item \textbf{改进}：
  \begin{itemize}
    \item RMQ 健康度上 Prometheus（\texttt{rabbitmq\_up} gauge）+ 告警规则。
    \item Cron 任务每次跑完写 \texttt{last\_run\_ts} 到 Redis，超过 10min 没更新就告警。
    \item DLX 队列堆积 > N 条告警（消费者 nack 进 DLX = 真有 bug）。
  \end{itemize}
\end{itemize}
\srcnote{cmd/main.go:49--58；initialize/cron.go:14--28 当前无心跳；repository/rabbitmq/order\_delay.go:98 nack requeue 后进 DLX}
\end{frame}

% ============================================================
\section{业务承诺总表}
% ============================================================

\begin{frame}{业务承诺总表 + 压测复核}
\begin{tabular}{@{}lll@{}}
\toprule
承诺 & 实现 & 复核 \\
\midrule
\textbf{30min 准点关单}    & RMQ TTL 30min                 & RMQ 单测 + 端到端 \\
\textbf{兜底关单}         & Cron 15min 窗口 / 5min 节奏    & RMQ 关停后 Cron 接管 \\
\textbf{幂等关单}         & SQL WHERE type=UnPaid          & 4 并发 → 1 单 cancel \\
\textbf{库存原子释放}     & DB tx 内 close + outbox        & Redis tx 外异步 release \\
\textbf{已付款绝不误关}   & WHERE 死锁 type=UnPaid         & 状态机不可回退 \\
\textbf{下游事件单次}     & outbox 在 tx 内写一次          & cancelled 事件唯一 \\
\bottomrule
\end{tabular}
\vspace{2mm}
\small 六条承诺缺一不可。压测复核：\texttt{orders/create} 单机 50{,}319 RPS / p95 2.33ms（远低于 P0 SLO 500ms）；幂等中间件 50 VU 15s 重放 755{,}033 次仍只 1 单；\textbf{RabbitMQ 未启动仍 0 错}——cmd/main.go:49--58 静默降级间接证明双保险的可降级性。
\srcnote{repository/rabbitmq/order\_delay.go + service/order\_cancel.go + service/order\_task.go 三件套；stressTest/REPORT.md \S 链路压测；cmd/main.go:49--58}
\end{frame}

\begin{frame}{回到那条 30 分钟的线}
\begin{itemize}
  \item 30 分钟不是延迟队列的 TTL 参数，是\textbf{商家、运营、风控、客服、财务}五方的合同。
  \item RMQ + Cron 不是技术冗余，是\textbf{业务承诺}：哪怕 RMQ 整个挂了，库存最多多锁 35min。
  \item CancelUnpaidOrder 幂等不是为了"代码优雅"，是因为\textbf{真的有 4 个调用方}会并发触发。
  \item 关单是金融事件，不是状态翻转——下游消费者必须设计成\textbf{不能给商家发货通知}。
  \item 三处遗留（口径 / orders/delete / cron 表达式）说明：\textbf{业务承诺到代码的最后一公里需要持续打磨}。
\end{itemize}
\keypoint{关单治理的核心不是"延迟队列怎么用"，是\textbf{把"30 分钟"这条业务线\\
落到 6 条承诺 + 双保险机制 + 幂等 SQL 上}。}
\end{frame}

% ============================================================
\appendix
% ============================================================

\begin{frame}{业务 Q\&A（上）}
\textbf{Q1}：为什么 30min，不能改成 1 小时？\\
\hspace{4mm}\small A：60min 的代价是双 11 锁单太久 / 爆款一直显示"已售罄" / 库存周转崩盘。30min 是商家投诉曲线和库存周转的工业界经验值。改要 5 个业务角色一起谈。

\vspace{2mm}
\textbf{Q2}：RMQ 挂了多久能发现？\\
\hspace{4mm}\small A：当前\textbf{没有外部告警}——cmd/main.go:49--58 只是 Warn 日志。改进项第 4 条。短期靠 Cron 兜底（15--20min 内库存释放）。

\vspace{2mm}
\textbf{Q3}：用户 APP 主动取消 + RMQ 同时到期会出错吗？\\
\hspace{4mm}\small A：不会。CloseOrderWithCheck 用 \texttt{WHERE type=UnPaid} 兜底——第一个成功的 RowsAffected=1，第二个 RowsAffected=0 直接 return。

\vspace{2mm}
\textbf{Q4}：客服后台关单和这条链路是同一个吗？\\
\hspace{4mm}\small A：业务侧应当统一走 \texttt{CancelUnpaidOrder}。如果客服有"强制关单"需求（用户付了但要退），那是\texttt{order.refunded} 链路，不在本 deck 范围内。
\end{frame}

\begin{frame}{业务 Q\&A（下）}
\textbf{Q5}：Cron 表达式 bug 为什么没人发现？\\
\hspace{4mm}\small A：因为幂等关单 SQL 兜住了功能正确性——CloseOrderWithCheck 看到已 Cancelled 直接 RowsAffected=0。但 DB QPS 被无谓放大 60 倍，监控里能看出 SELECT type=UnPaid 异常高频。改成 \texttt{0 */5 * * * *} 即可。

\vspace{2mm}
\textbf{Q6}：为什么不直接用 Redis ZSet 做延迟队列，省掉 RMQ？\\
\hspace{4mm}\small A：service/order.go:88--92 确实有 ZSet 但只用作内部状态。换成 ZSet 替代 RMQ 也行，但要自己写消费循环 + ack + 重试 + DLX——RMQ 把这些都打包了。当前业务量下 RMQ 性价比更高。

\vspace{2mm}
\textbf{Q7}：对账脚本什么时候上？\\
\hspace{4mm}\small A：运维 backlog 第 1 项。当前依赖客服反馈兜底——商家投诉"库存对不上"才去查 GetStockSnapshot。不是好状态，需要补凌晨 cron。

\vspace{2mm}
\textbf{Q8}：双保险设计能保证 100\% 不漏单吗？\\
\hspace{4mm}\small A：不能。还有一条漏：DB tx 提交了 outbox 但 outbox publisher 挂了 → 下游 ES / 客服系统不知道关单。这就是为什么需要\textbf{对账}作为最终兜底。
\end{frame}

\begin{frame}{代码位置一览}
\begin{description}
  \item[30min 常量]\srcnote{repository/rabbitmq/order\_delay.go:21--22}
  \item[RMQ 延迟拓扑]\srcnote{repository/rabbitmq/order\_delay.go:24--56}
  \item[RMQ publish]\srcnote{repository/rabbitmq/order\_delay.go:58--72}
  \item[RMQ consume]\srcnote{repository/rabbitmq/order\_delay.go:74--105}
  \item[CancelUnpaidOrder]\srcnote{service/order\_cancel.go:19--63}
  \item[CloseOrderWithCheck 幂等 SQL]\srcnote{repository/db/dao/order.go:185--196}
  \item[Cron 任务]\srcnote{service/order\_task.go:14--48}
  \item[Cron 注册 (含表达式 bug)]\srcnote{initialize/cron.go:13--28}
  \item[RMQ 启动失败静默降级]\srcnote{cmd/main.go:49--58}
  \item[下单时发 RMQ + ZSet]\srcnote{service/order.go:88--98}
  \item[orders/delete 绕过关单链路]\srcnote{routes/router.go:81；service/order.go:168--192}
  \item[Cron 兜底窗口 15min]\srcnote{service/order\_task.go:16；repository/db/dao/order.go:162--170}
\end{description}
\end{frame}

\end{document}
